\chapter{Diseño del sistema y marco de evaluación}
\label{cap:evaluacion}

Este capítulo presenta una visión general del pipeline y, sobre todo, el marco de evaluación común a todos los módulos: qué métricas se usan y cómo interpretarlas, cómo se valida sin fuga de datos y qué pruebas estadísticas respaldan cada comparación. Reunir estas definiciones aquí evita repetirlas en cada módulo y permite que los capítulos siguientes se centren en sus resultados.

\section{Visión general del pipeline}

El sistema se organiza en un pipeline secuencial que transforma los datos brutos de campaña en predicciones y recomendaciones accionables, agrupado en cuatro bloques funcionales:
\begin{itemize}
    \item \emph{Bloque I — Preparación de datos} (\texttt{01\_clean}, \texttt{02\_geo}, \texttt{03\_users}): limpieza, construcción de tablas dimensionales y enriquecimiento demográfico (Capítulo~\ref{cap:datos}).
    \item \emph{Bloque II — Segmentación} (\texttt{05\_demo\_segmentation}, \texttt{06\_segmentation}): dos niveles de agrupación de usuarios, M0 (demográfica) y M1 (comportamental) (Capítulo~\ref{cap:segmentacion}).
    \item \emph{Bloque III — Modelado predictivo} (\texttt{07\_propensity}, \texttt{08\_recommender}): propensión al clic (M2, Capítulo~\ref{cap:propension}) y recomendador por cluster (M3, Capítulo~\ref{cap:recomendador}).
    \item \emph{Bloque IV — Previsión temporal} (\texttt{09\_forecast}): predicción del volumen semanal de clics (M4, Capítulo~\ref{cap:forecast}).
\end{itemize}
Como análisis complementarios se incluyen un diagnóstico de estabilidad del número de clusters (\texttt{05b\_k\_stability}), un modelo inverso producto\,$\rightarrow$\,perfil (\texttt{10\_inverse\_profile}) y un ejemplo de despliegue con análisis de retorno (Capítulo~\ref{cap:negocio}). Cada capítulo de módulo presenta de forma conjunta su \emph{método} y sus \emph{resultados}, para facilitar la lectura.

\section{Métricas de evaluación}
\label{sec:metricas}

Cada tipo de problema exige métricas distintas. A continuación se definen las empleadas y, lo más importante, cómo interpretarlas: qué valor es bueno y cuál es el de referencia o ``azar''.

\subsection{Propensión al clic (clasificación binaria)}

El problema es de clases desbalanceadas (pocos clics frente a muchas no-interacciones), lo que condiciona la elección de métricas: la \emph{exactitud} (\textit{accuracy}) engaña, porque un modelo que prediga ``nunca clic'' acertaría la mayoría de las veces sin ser útil.

\begin{itemize}
    \item \textbf{PR-AUC} (área bajo la curva \textit{precision-recall}), métrica principal. Resume cómo de bien el modelo identifica la clase rara (el clic) a distintos umbrales. Va de 0 a 1; mayor es mejor. Su valor de \emph{referencia} (lo que daría el azar) no es 0,5 sino la prevalencia de positivos: aquí $\approx0{,}238$. Por eso un PR-AUC de 0,40 sobre una prevalencia de 0,24 sí indica señal real. Se prefiere a la AUC-ROC en desbalanceo porque se centra en la clase positiva~\cite{davisgoadrich2006,saito2015}.
    \item \textbf{AUC-ROC} (área bajo la curva ROC)~\cite{fawcett2006}, métrica de apoyo. Mide la capacidad de \emph{ordenar} un positivo por encima de un negativo elegidos al azar. Va de 0 a 1: 0,5 es azar, 1 es perfecto. Orientativamente, $\approx0{,}7$ es modesto, $0{,}8$ bueno y $\geq0{,}9$ excelente. En desbalanceo tiende a ser \emph{optimista}, de ahí que la principal sea PR-AUC.
    \item \textbf{Brier score}~\cite{brier1950}, calidad de la \emph{calibración}. Mide si las probabilidades son \emph{realistas} (no solo si el orden es correcto): es el error cuadrático medio entre la probabilidad predicha y el resultado real (0/1). Menor es mejor (0 = perfecto).
    \item \textbf{F1, MCC y \textit{Balanced Accuracy}}, calidad de la decisión binaria a un umbral fijo (aquí, la prevalencia 0,238). El F1 es la media armónica de precisión y exhaustividad (0 a 1). El MCC (coeficiente de correlación de Matthews~\cite{matthews1975}) resume la matriz de confusión completa y es robusto al desbalanceo: va de $-1$ a $1$, con 0 = azar. La \textit{Balanced Accuracy} promedia el acierto en cada clase por separado (0,5 = azar).
\end{itemize}

\subsection{Recomendación (calidad del \textit{ranking})}

Aquí no se predice un número, sino un \emph{orden} de categorías; las métricas miden si las categorías relevantes para el usuario aparecen arriba en una lista de tamaño $K$.
\begin{itemize}
    \item \textbf{Hit Rate@$K$} (HR@$K$): fracción de usuarios para los que \emph{al menos una} categoría relevante está entre las $K$ recomendadas. Responde ``¿acerté algo en el top-$K$?''. Va de 0 a 1.
    \item \textbf{NDCG@$K$} (\textit{Normalized Discounted Cumulative Gain})~\cite{jarvelin2002}: como HR pero premia el orden, dando más peso a los aciertos en las primeras posiciones. De 0 a 1.
    \item MAP@$K$ (\textit{Mean Average Precision}): precisión media a lo largo de la lista, también sensible al orden. De 0 a 1.
\end{itemize}

\subsection{Previsión temporal}

Para el volumen semanal de clics se usan métricas de error de predicción (menor es mejor):
\begin{itemize}
    \item \textbf{MAE} (error absoluto medio) y \textbf{RMSE} (raíz del error cuadrático medio): en las mismas unidades que la serie (clics/semana); el RMSE penaliza más los errores grandes.
    \item MAPE (error porcentual absoluto medio): el error relativo en \%; cómodo de interpretar pero inestable cuando los valores reales son pequeños.
\end{itemize}

\section{Validación sin fuga de datos}
\label{sec:validacion}

La estimación fiable del rendimiento exige evitar la \emph{fuga de datos} (\textit{data leakage}): que el modelo acceda en entrenamiento a información que no tendría en producción, lo que infla artificialmente las métricas~\cite{kaufman2012leakage}. Se emplean tres esquemas según el módulo:
\begin{itemize}
    \item \emph{Validación cruzada agrupada por usuario} (\texttt{GroupKFold}, M2): reparte los datos en 5 \emph{folds} garantizando que ningún usuario esté a la vez en entrenamiento y validación. Es imprescindible aquí porque las variables de un usuario son idénticas en todos sus eventos; sin agrupar, el modelo ``memorizaría'' al usuario y sobreestimaría su capacidad. Mide la generalización a usuarios \emph{nuevos}.
    \item \emph{Validación cruzada temporal} (\texttt{TimeSeriesSplit}, modelo \textit{warm} y M4): el entrenamiento es siempre \emph{pasado} y la validación \emph{futuro}, respetando el orden cronológico.
    \item \emph{Reconstrucción solo con \textit{train}} (M3): el segmento de cada usuario se recalcula usando únicamente los datos de entrenamiento, para que no ``conozca'' el periodo de test.
\end{itemize}

\paragraph{Corrección de prior (\textit{prior shift}).} El modelo de propensión se entrena con una muestra \emph{balanceada} (\textasciitilde24\,\% de clics) pero en producción la tasa real es $\approx 2$\,\%; sin ajuste, las probabilidades saldrían infladas. La transformación de Dal Pozzolo et al.~\cite{dalpozzolo2015} (ecuación~\ref{eq:prior_correction}) las reescala al prior real. Es monótona: cambia la magnitud de las probabilidades pero \emph{no} su orden, por lo que no afecta a las métricas de \textit{ranking} (PR-AUC, AUC-ROC).

\paragraph{Validación cruzada frente a \textit{bootstrap}.} Son complementarios y responden a preguntas distintas. La validación cruzada mide \emph{cómo de bien generaliza} el modelo (reentrena $k$ veces, cada dato es test una vez). El \textit{bootstrap} mide la \emph{incertidumbre} de una métrica ya calculada: remuestrea el conjunto de test al azar muchas veces (con reemplazo) y observa cuánto varía la métrica, lo que da un intervalo de confianza sin necesidad de reentrenar.

\section{Contraste estadístico}
\label{sec:tests}

Que un modelo obtenga una métrica algo mejor que otro \emph{no} basta para afirmar que es superior: la diferencia podría deberse al azar de la muestra concreta. Las pruebas estadísticas cuantifican esa duda. La idea común es el \emph{$p$-valor}: la probabilidad de observar una diferencia tan grande como la medida \emph{si en realidad no hubiera diferencia}; un $p<0{,}05$ se interpreta como evidencia de que la diferencia no es casualidad. El uso de pruebas no paramétricas para comparar modelos sigue las recomendaciones de Dem\v{s}ar~\cite{demsar2006}. Las pruebas empleadas son:

\begin{itemize}
    \item \textbf{Intervalo de confianza \textit{bootstrap}}~\cite{efron1979}: se remuestrea el test con reemplazo (1.000 veces) y se recalcula la métrica o su diferencia. Si el intervalo al 95\,\% de una diferencia \emph{no incluye el 0}, la diferencia es robusta. No depende del número de \emph{folds}.
    \item \textbf{Test de Wilcoxon (pareado)}~\cite{wilcoxon1945}: compara dos modelos \emph{fold a fold}, sin asumir normalidad. \emph{Limitación} importante: con 5 \emph{folds} el menor $p$-valor posible es 0,0625, así que un resultado no significativo puede deberse a falta de potencia, no de diferencia; por eso se apoya en el \textit{bootstrap}.
    \item \textbf{Test de Friedman}~\cite{friedman1937}: prueba \emph{global} para comparar tres o más métodos sobre los mismos sujetos (``¿hay alguna diferencia entre ellos?''). Se hace antes de comparar pares.
    \item \textbf{Test de McNemar}~\cite{mcnemar1947}: compara dos métodos \emph{caso a caso} en acierto/fallo, fijándose solo en los casos en que discrepan. Útil para contrastar la tasa de acierto (HR@5) entre el recomendador y un \textit{baseline}.
    \item \textbf{Corrección de Holm}~\cite{holm1979}: al realizar \emph{varias} comparaciones a la vez, la probabilidad de un falso positivo se acumula. Holm endurece el umbral de significación de forma secuencial (al $p$ más pequeño le exige el listón más estricto), controlando ese riesgo.
    \item \textbf{Test de Diebold-Mariano (DM)}~\cite{diebold1995}: el equivalente para comparar \emph{predicciones de series temporales}; contrasta los errores semana a semana de dos modelos de \textit{forecast}. Con muestras pequeñas se aplica la corrección de Harvey-Leybourne-Newbold~\cite{harvey1997} (factor de muestra reducida y $t$ de Student), imprescindible aquí por disponer de pocos orígenes.
    \item \emph{Tamaños de efecto}: el $p$-valor dice si una diferencia existe, no \emph{cómo de grande} es. Se reportan la W de Kendall (para el Friedman; 0 = sin efecto, 1 = efecto máximo) y la correlación rango-biserial $|r|$ (para los Wilcoxon pareados; orientativamente, $0{,}1$ pequeño, $0{,}3$ medio, $0{,}5$ grande).
\end{itemize}

Todas las semillas aleatorias se fijan en 42 para garantizar la reproducibilidad (Apéndice~\ref{ap:reproducibilidad}).
